\documentclass[11pt,a4paper]{article}
\usepackage[T1]{fontenc}
\usepackage{lmodern,amsmath,amssymb,amsthm,mathtools,microtype}
\usepackage[margin=26mm]{geometry}
\usepackage[colorlinks=true,linkcolor=blue,citecolor=blue,urlcolor=blue]{hyperref}
\newtheorem{theorem}{Theorem}[section]
\newtheorem{proposition}[theorem]{Proposition}
\newtheorem{lemma}[theorem]{Lemma}
\newtheorem{corollary}[theorem]{Corollary}
\theoremstyle{definition}
\newtheorem{definition}[theorem]{Definition}
\newtheorem{remark}[theorem]{Remark}
\newcommand{\Q}{\mathbb Q}
\newcommand{\Z}{\mathbb Z}
\newcommand{\C}{\mathbb C}
\newcommand{\hF}{h_{\mathrm F}}
\newcommand{\OK}{\mathcal O_K}
\newcommand{\rad}{\operatorname{rad}}
\newcommand{\Jac}{\operatorname{Jac}}
\newcommand{\Ardeg}{\widehat{\deg}}
\setlength{\emergencystretch}{2em}
\title{Evaluation of Differentials and Heights of Modular Jacobians}
\author{}
\date{9 November 2026}
\begin{document}
\maketitle
\begin{abstract}
We bound the stable Faltings height of a Jacobian by the discriminant
of one unramified fibre of a Belyi map. The proof evaluates integral
regular differentials in a trace-dual lattice and estimates the
complex evaluation maps by the mean-value inequality. For a modular
curve of degree $d$ over the $j$-line and level $N$, the different of
the elliptic torsion scheme gives an explicit bound of order
$d\log(2N)$, uniformly in the subgroup defining the curve.
Applied to the modular realization
underlying the five-contact Diophantine construction, it yields
\[
 h(a)<2^{330770}M_5(a)^3(1+\log M_5(a))
       <2^{330789}M_5(a)^{7/2}.
\]
The height $h$ is logarithmic. The geometric and modular realization
is stated as the arithmetic input; no additional $abc$ or Szpiro
hypothesis is imposed. All bounds use the original stable Faltings
metric, and the numerical constants admit exact verification.
\end{abstract}

\section{Conventions and the evaluation estimate}

If $K$ is a number field, $\OK$ denotes its ring of integers.
For a nonzero ideal $\mathfrak a\subset\OK$, write
$\mathrm N\mathfrak a=\#(\OK/\mathfrak a)$.
If $E$ is a projective $\OK$-module with Hermitian metrics at
its complex embeddings, $\Ardeg(E)$ means the arithmetic degree
of its determinant. In rank one, for a nonzero integral section $s$,
\[
 \Ardeg(E)=\log\#(E/\OK s)
                  -\sum_{\sigma:K\hookrightarrow\C}\log\|s\|_\sigma.
\]
The definition is independent of the choice of $s$ by the product
formula. Both embeddings in a complex conjugate pair are included.

Throughout, $\hF$ is the original stable Faltings height.
For a curve $X/K$ with semistable reduction, its Jacobian height is
\begin{equation}\label{eq:height-convention}
 \hF(\Jac X)=\frac{\Ardeg H^0(\mathcal X,
                                  \omega_{\mathcal X/\OK})}{[K:\Q]},
\end{equation}
where $\mathcal X$ is a minimal regular semistable model. The metric
on differentials at a complex embedding is
\[
 \langle\omega,\eta\rangle
       =\frac i2\int_{X(\C)}\omega\wedge\overline\eta.
\]
This is the determinant metric corresponding to the original Faltings
normalization; see \cite[Sections~2.1 and~2.3]{Javanpeykar} and
\cite[Section~2.1]{vonKaenel}. Passing to a finite extension on which
the curve is semistable does not change the stable height.

A \emph{Belyi map} is a finite morphism from a smooth projective
geometrically connected curve to $\mathbb P^1$, unramified outside
$0,1,\infty$. Let
\[
 \pi:X\longrightarrow\mathbb P^1
\]
be such a map, of degree $d$, where $X$ has genus $g\ge1$.
Write $t$ for the coordinate on the target.
Choose a number field $K$ over which $X$ and $\pi$ are defined,
$X$ is semistable, and every point of $\pi^{-1}(2)$ is rational.
Let $\mathcal Y$ be the normalization of $\mathbb P^1_{\OK}$ in
$K(X)$. It is finite and flat over $\mathbb P^1_{\OK}$.
On the affine $t$-chart write its algebra as $B$, a finite projective
$\OK[t]$-algebra of rank $d$, and put
\begin{equation}\label{eq:order}
 C=B/(t-2)B\subset K^d,
 \qquad O=\OK^d.
\end{equation}
The inclusion in $K^d$ uses the $d$ distinct points of the split
fibre. Every element of $C$ is integral over $\OK$, so $C\subset O$.
Let $\mathfrak D(C)$ be the discriminant ideal of its trace pairing.

\begin{theorem}\label{thm:evaluation}
With this notation,
\begin{equation}\label{eq:evaluation-bound}
 \hF(\Jac X)
 \le\frac{\log\mathrm N\mathfrak D(C)}{2[K:\Q]}
       +g\log\frac2{\sqrt\pi}.
\end{equation}
\end{theorem}

The auxiliary field $K$ is used to split a finite fibre and to
compute a stable height. Its degree will cancel from the estimates.
No upper bound for $[K:\Q]$ is assumed in this theorem.

\section{Trace-dual lattices and regular differentials}

\begin{lemma}\label{lem:trace-lattice}
Let $C\subset O=\OK^d$ be an order with $C\otimes K=K^d$, and
let
\[
 C^\vee=\{x\in K^d:\operatorname{Tr}_{K^d/K}(xC)\subset\OK\}.
\]
For any set $I$ of $g$ coordinate positions, set
$F=\operatorname{pr}_I(C^\vee)\subset K^g$.
Then $F$ is a projective $\OK$-module containing $\OK^g$.
With the standard Euclidean metric on $K^g$ at every embedding,
\[
 \Ardeg(F)\le\tfrac12\log\mathrm N\mathfrak D(C).
\]
\end{lemma}
\begin{proof}
Work locally at a finite prime $\mathfrak p$. The trace pairing
on $K^d$ is the ordinary dot product, and $O$ is self-dual.
Thus $O\subset C^\vee$, and dualizing the inclusion $C\subset O$
shows
\[
 \operatorname{length}(C^\vee/O)
 =\operatorname{length}(O/C)
 =\tfrac12v_{\mathfrak p}(\mathfrak D(C)).
\]
For example, if a basis matrix for $C$ relative to $O$ is $U$,
then the trace Gram matrix is $U^{\mathsf T}U$, which proves the
last equality. These equalities are valid locally even when
$C$ is not globally free.

Projection gives a surjection
$C^\vee/O\twoheadrightarrow F/\OK^g$. The latter quotient
therefore has no larger length at any prime.
A finite torsion-free module over the Dedekind ring $\OK$ is
projective. Finally, the standard basis of $\OK^g$ has determinant
norm one, so
\[
 \Ardeg(F)=\log\#(F/\OK^g)
 \le\tfrac12\log\mathrm N\mathfrak D(C).
\]
\end{proof}

\begin{lemma}\label{lem:integral-evaluation}
Let $\mathcal X$ be a minimal regular semistable model of $X$.
The evaluations of a differential in
$E=H^0(\mathcal X,\omega_{\mathcal X/\OK})$ at the points of
$\pi^{-1}(2)$, expressed in the coordinate $t$, lie in $C^\vee$.
Equivalently, there is an integral map
\[
 E\longrightarrow C^\vee,
 \qquad \omega\longmapsto
       \left(\frac{\omega}{dt}(P)\right)_{P\in\pi^{-1}(2)}.
\]
\end{lemma}
\begin{proof}
First, normalization is finite because the base is excellent.
A normal surface is Cohen--Macaulay, and a finite morphism from
such a surface to a regular surface is flat by miracle flatness;
see \cite[Tag~00R4]{Stacks}. Consequently $B$ is finite projective
over $\OK[t]$.

Take a common regular model $\mathcal Z$ dominating
$\mathcal X$ and $\mathcal Y$. Pullback under a blowup of regular
arithmetic surfaces takes regular relative differentials to regular
relative differentials. One can also read this directly from the
canonical-sheaf formula for a blowup; see
\cite[Lemma~2.3.1]{Javanpeykar}.
Hence a differential in $E$ is regular on $\mathcal Z$.
The birational morphism $\mathcal Z\to\mathcal Y$ is an
isomorphism at every codimension-one point of the normal surface
$\mathcal Y$.

The dualizing module of a Cohen--Macaulay normal surface is
reflexive: it is torsion-free and satisfies Serre's condition
$(S_2)$. Thus sections defined at all codimension-one points extend
across the codimension-two complement; see
\cite[Tags~0AWQ and~0AVT, Lemmas~31.13.12--31.13.14]{Stacks}.
The differential therefore defines a section of the relative
dualizing sheaf of $\mathcal Y$.

Finite duality identifies this sheaf on the affine chart with
\[
 \operatorname{Hom}_{\OK[t]}(B,\OK[t])\,dt;
\]
see \cite[Tags~0BUL and~0C0I]{Stacks}.
This statement uses the dualizing \emph{module}; it does not require
that the normal surface have an invertible dualizing sheaf.
Because $B$ is finite projective, specializing the dual at $t=2$
gives $\operatorname{Hom}_{\OK}(C,\OK)\,dt$.
On the generic fibre the identification sends a differential
$f\,dt$ to the functional $b\mapsto\operatorname{Tr}(fb)$.
After splitting the fibre, its coordinates are exactly the values
$f(P)=\omega/dt(P)$. This proves the assertion.
\end{proof}

\begin{proof}[Proof of Theorem~\ref{thm:evaluation}]
Riemann--Hurwitz gives
\[
 2g-2=-2d+\deg\operatorname{Ram}(\pi)\le d-3,
\]
since each of the three branch fibres contributes at most $d-1$.
In particular $d>2g-2$. A nonzero holomorphic differential cannot
vanish at all $d$ distinct points of $\pi^{-1}(2)$. Therefore the
evaluation map has rank $g$. Choose $g$ coordinate positions $I$
for which projection remains injective. By
Lemmas~\ref{lem:trace-lattice} and~\ref{lem:integral-evaluation},
we obtain an integral, generically invertible map
\[
 \phi:E\longrightarrow F=\operatorname{pr}_I(C^\vee),
 \qquad\Ardeg(F)\le\tfrac12\log\mathrm N\mathfrak D(C).
\]

Fix a complex embedding of $K$. The disc
$D=\{t:|t-2|<1/2\}$ is simply connected and avoids the branch
locus. Its inverse image is a disjoint union of $d$ analytic
sheets, each mapped isomorphically to $D$. On one such sheet write
$\omega=f(t)\,dt$. The mean-value inequality implies
\[
 |f(2)|^2\le\frac4\pi\int_D|f(t)|^2\,dx\,dy
       \le\frac4\pi\|\omega\|_{L^2}^2.
\]
Hence the norm of each evaluation functional is at most
$2/\sqrt\pi$. In an orthonormal basis of $E\otimes_\sigma\C$,
each row of the evaluation matrix has this same bound on its
Euclidean norm. Hadamard's determinant inequality gives
\[
 \|\det\phi\|_\sigma\le(2/\sqrt\pi)^g.
\]
The determinant is nonzero and integral. Taking its arithmetic
degree yields
\[
 \Ardeg(E)\le\Ardeg(F)
       +\sum_\sigma\log\|\det\phi\|_\sigma
 \le\tfrac12\log\mathrm N\mathfrak D(C)
       +[K:\Q]g\log(2/\sqrt\pi).
\]
Now apply \eqref{eq:height-convention}.
\end{proof}

\section{The discriminant of the fibre}

Let $S$ be a finite set of rational primes, and put
\[
 R_S=\prod_{p\in S}p,\qquad \omega(S)=\#S.
\]
Assume that the normalization defining the Belyi cover is finite
\'etale over
\begin{equation}\label{eq:good-open}
 \operatorname{Spec}\mathcal O_{K_0}
       [S^{-1},t,t^{-1},(t-1)^{-1}],
\end{equation}
where $K_0$ is a field of definition.
This property persists after any finite field extension:
finite \'etale base change is finite \'etale, and a finite \'etale
algebra over a normal base is normal. Thus it agrees with the
normalization of the new generic cover on this open set.
In particular the auxiliary splitting field in Section~1 introduces
no new vertical branch primes outside $S$.

\begin{lemma}\label{lem:discriminant}
Under \eqref{eq:good-open}, the order $C$ in \eqref{eq:order}
satisfies
\[
 \frac{\log\mathrm N\mathfrak D(C)}{[K:\Q]}
 \le d\left(\log2+\log R_S+\omega(S)\log d\right).
\]
\end{lemma}
\begin{proof}
The determinant of the trace pairing of $B/\OK[t]$ defines an
effective discriminant divisor on the base. Formation of the trace
pairing commutes with specialization. Its intersection with the
section $t=2$ is therefore $\log\mathrm N\mathfrak D(C)$.
The coefficient of this divisor at a codimension-one point is
$\sum_j f_jr_j$, where $e_j,f_j,r_j$ are the ramification index,
residue degree, and different exponent of the corresponding
extensions of discrete valuation rings. They satisfy
$\sum_j e_jf_j=d$. These are the usual local
trace-discriminant and different identities; see
\cite[Chapter~III, Sections~3 and~6]{Serre} and
\cite[Section~4.2]{Javanpeykar}.

At a horizontal branch divisor the residue characteristic is zero.
The different exponent is $r_j=e_j-1$, so the discriminant
coefficient is at most $d$. The section $t=2$ meets $t=0$
with normalized intersection $\log2$, and meets neither $t=1$
nor $t=\infty$. Thus the horizontal contribution is at most
$d\log2$.

Let $\mathfrak p\mid p$ be a prime of $K$, and write
$e_{\mathfrak p}=v_{\mathfrak p}(p)$.
Lenstra's different bound, in the form
\cite[Proposition~4.1.3]{Javanpeykar}, gives
\[
 r_j\le e_j-1+e_j e_{\mathfrak p}
                \left\lfloor\frac{\log d}{\log p}\right\rfloor.
\]
This bound also applies to the possibly imperfect residue field
at the generic point of a vertical fibre.
Its discriminant coefficient is consequently at most
\[
 d\left(1+e_{\mathfrak p}
                \left\lfloor\frac{\log d}{\log p}\right\rfloor\right).
\]
Only $p\in S$ can contribute. The intersection of the vertical
fibre over $\mathfrak p$ with the section has value
$f_{\mathfrak p}\log p$. Using
\[
 \sum_{\mathfrak p\mid p}e_{\mathfrak p}f_{\mathfrak p}=[K:\Q],
 \qquad \sum_{\mathfrak p\mid p}f_{\mathfrak p}\le[K:\Q],
\]
and
$\lfloor\log d/\log p\rfloor\log p\le\log d$
proves the estimate.
\end{proof}

\begin{theorem}\label{thm:coverheight}
For a Belyi cover satisfying \eqref{eq:good-open},
\begin{equation}\label{eq:coverheight}
 \hF(\Jac X)
 \le\frac d2\left(\log2+\log R_S+\omega(S)\log d\right)
       +g\log\frac2{\sqrt\pi}.
\end{equation}
\end{theorem}
\begin{proof}
Combine Theorem~\ref{thm:evaluation} and
Lemma~\ref{lem:discriminant}.
\end{proof}

\section{The integral modular cover}

Let $N$ be a positive integer and let
$H\subset(\Z/N\Z)^\times$ contain $-1$. Write
\[
 X_H=X_1(N)/H,\qquad t_H=[(\Z/N\Z)^\times:H],
 \qquad J_H=\Jac X_H.
\]
The action is the diamond action. Assume that $X_H$ has positive
genus, and use the rational map $t=j/1728$ to the projective line.

\begin{lemma}\label{lem:modular-model}
The map $X_H\to\mathbb P^1_t$ has degree
\[
 d_H=Nt_H\prod_{p\mid N}(1+p^{-1}).
\]
It is unramified outside $0,1,\infty$ and admits a finite \'etale
model over
\[
 U=\operatorname{Spec}\Z[1/(6N),t,t^{-1},(t-1)^{-1}].
\]
\end{lemma}
\begin{proof}
Over $U$ consider
\begin{equation}\label{eq:universal-curve}
 E_t:\quad y^2=x^3-\frac{3t}{t-1}x-\frac{2t}{t-1}.
\end{equation}
Its discriminant and $j$-invariant are
\[
 \Delta(E_t)=\frac{1728t^2}{(t-1)^3},\qquad j(E_t)=1728t.
\]
Thus it is an elliptic curve over $U$. Since $N$ is invertible,
its scheme of points of exact order $N$ is finite \'etale.
The constant group $H$ acts freely on this scheme: if $aP=P$
and $P$ has exact order $N$, then $a=1$ in $(\Z/N\Z)^\times$.
The quotient is finite \'etale. This follows also from the
classification of finite \'etale covers by locally constant finite
sets; freeness is enough, even when a residue characteristic divides
$|H|$.

To identify its generic fibre, fix an elliptic curve with
$j$-invariant different from $0$ and $1728$ over an algebraic
closure. Its automorphism group is $\{1,-1\}$. The geometric
fibre of $X_1(N)$ is the set of primitive torsion points modulo
this automorphism group. Since $-1\in H$, the fibre of $X_H$
is precisely the set of primitive torsion points modulo $H$.
Different isomorphisms to \eqref{eq:universal-curve} differ by
$\pm1$ and hence give the same $H$-orbit. In particular quadratic
twisting introduces no extra double cover. This identifies the
finite \'etale model with the normalization of $X_H$ over $U$.

There are $N^2\prod_{p\mid N}(1-p^{-2})$ primitive points.
Division by $|H|=\varphi(N)/t_H$ gives the asserted degree.
Equivalently, the congruence subgroup contains $-I$, so its
$\operatorname{SL}_2(\Z)$ index is already the degree of the
coarse map to $X(1)$; see \cite[Section~4]{vonKaenel}.
The only possible branch points of that map are the two elliptic
points and the cusp, which become $0,1,\infty$ in the $t$ coordinate.
\end{proof}

\begin{lemma}\label{lem:torsion-different}
Let $V$ be an excellent discrete valuation ring of characteristic zero,
let $E/V$ be an elliptic scheme, and let $n$ be a positive integer.
Write $E[n]=\operatorname{Spec} A$, and let $\widetilde A$ be its
normalization in its generic finite \'etale algebra. Under the generic
trace pairing, the complementary modules satisfy
\[
 A^\vee=n^{-1}A,\qquad
 \widetilde A^\vee\subseteq n^{-1}\widetilde A.
\]
The different of every field component of $\widetilde A/V$ therefore
contains $n$. The same assertion holds for every intermediate field
of such a component.
\end{lemma}
\begin{proof}
Multiplication $[n]:E\to E$ is finite flat and a local complete
intersection morphism. Its derivative in invariant differential bases
is multiplication by $n$. Consequently its relative differential
module has zeroth Fitting ideal $(n)$. For a finite flat local complete
intersection, this K\"ahler different is the different; see
\cite[Tag~0BWB, Lemmas~49.12.3 and~49.12.6]{Stacks}.

Base change along the zero section gives $E[n]/V$. The relative
differential module and its zeroth Fitting ideal commute with base
change, while the morphism remains finite flat and a local complete
intersection. Its different is therefore $(n)$. Its relative
dualizing module is invertible. Under the trace pairing this is the
complementary fractional ideal $A^\vee$, whose inverse is the different;
see \cite[Tag~0BW0, Lemmas~49.8.2--49.8.3, and Tag~0BWB,
Lemma~49.12.4]{Stacks}. Thus $A^\vee=n^{-1}A$ as fractional ideals,
not merely as abstract modules.

Normalization is finite, and inclusion reverses trace duals:
\[
 \widetilde A^\vee\subseteq A^\vee=n^{-1}A
                              \subseteq n^{-1}\widetilde A.
\]
Hence $n\widetilde A^\vee\subseteq\widetilde A$, which says exactly
that the different contains $n$. Finally, for normal discrete
valuation rings $V\subseteq W\subseteq U$, the different formula
\[
 \mathfrak D_{U/V}=\mathfrak D_{U/W}\mathfrak D_{W/V}U
\]
shows that $v_U(\mathfrak D_{U/V})\le v_U(n)$ implies
$v_W(\mathfrak D_{W/V})\le v_W(n)$, because the different exponent
of $U/W$ is nonnegative. Apply this in each field component and
at each place. These finite-duality and different arguments do not
require a perfect residue field.
\end{proof}

\begin{corollary}\label{cor:modularheight}
Let $g_H$ be the genus of $X_H$. Then
\begin{equation}\label{eq:modularheight}
 \hF(J_H)\le
 \frac{d_H}{2}\left(\log12+\log N+2\log d_H\right)
                  +g_H\log\frac2{\sqrt\pi}.
\end{equation}
\end{corollary}
\begin{proof}
Use the auxiliary field $K$ from Theorem~\ref{thm:evaluation}.
At a prime $\mathfrak p$ of $K$ over $p>3$, localize the $t$-line
at the generic point of the special fibre. The family
\eqref{eq:universal-curve} is an elliptic scheme over this discrete
valuation ring: its discriminant $1728t^2/(t-1)^3$ is a unit there.
By Lemma~\ref{lem:modular-model}, the generic cover for $X_H$ is
the primitive $N$-torsion cover modulo $H$. Each of its field
components is an intermediate field of a field generated by a single
$N$-torsion point. Lemma~\ref{lem:torsion-different} therefore bounds
every different exponent above $\mathfrak p$ by the valuation of $N$.
The coefficient of the vertical discriminant divisor is at most
\[
                    d_H v_{\mathfrak p}(N).
\]
After summing over these primes and dividing by $[K:\Q]$, the
contribution is at most $d_H\log N$.

At $p=2,3$, the estimate in Lemma~\ref{lem:discriminant} gives at
most $d_H(\log6+2\log d_H)$. The horizontal intersection at
$t=2$ contributes at most $d_H\log2$. Thus the normalized
logarithmic discriminant of the fibre order is at most
\[
                 d_H(\log12+\log N+2\log d_H).
\]
Apply Theorem~\ref{thm:evaluation}. In particular no degree of a
full torsion splitting field enters this estimate.
\end{proof}

\section{An abelian surface occurring as an isogeny factor}

\begin{lemma}\label{lem:surface}
Let $J$ be an abelian variety of dimension $g\ge2$ over a number
field $L'$, and let $A$ be an abelian surface which is an isogeny
factor of $J$ over $L'$. Put $d'=[L':\Q]$ and
\[
 H_J=\max(1,\hF(J)+3g/2).
\]
Then
\begin{equation}\label{eq:surface}
 \hF(A)\le H_J+
 8\left(32\log14+\log d'+
                \log\max(1,\log d',H_J)\right).
\end{equation}
\end{lemma}
\begin{proof}
Choose an $L'$-homomorphism $A\to J$ with finite kernel, and
let $A_0\subset J$ be its image. It is an abelian surface over
$L'$ and is $L'$-isogenous to $A$.
The stable version of R\'emond's subvariety inequality
\cite[Corollary~1.2 and the paragraph following Corollary~1.3]{Remond}
gives
\[
 \hF(A_0)+\hF(J/A_0)\le\hF(J).
\]
Bost's lower bound in the original normalization is
$\hF(B)\ge-\log(\pi\sqrt2)\dim B>-3\dim B/2$.
Thus $\hF(A_0)\le H_J$.

Apply \cite[Theorem~1.9(1)]{GaudronRemond} with source $A_0$
and target $A$, both over $L'$. In dimension two it supplies an
isogeny of degree at most
\[
 \left(14^{32}d'\max(1,\log d',\hF(A_0))\right)^{16}.
\]
Stable Faltings height increases under an isogeny of degree $D$
by at most $\frac12\log D$; see also
\cite[Proposition~2.1]{Remond}. Replacing $\hF(A_0)$ by $H_J$
proves \eqref{eq:surface}.
\end{proof}

The field $L'$ in this lemma is the field on which the surface
factor is realized. It need not split a fibre of a modular curve,
and it need not be the auxiliary field $K$ of
Theorem~\ref{thm:evaluation}. The evaluation estimate is a stable
height bound valid after arbitrary extension; the degree entering
\eqref{eq:surface} is only $[L':\Q]$.

\section{The five-contact height estimate}

For $a=u/v\in\Q$ in lowest terms, with $v>0$, put
$h(a)=\log\max(|u|,v)$. Define
\[
 S_*=\{p\text{ prime}:p\le76176\},\quad
 B_*=\prod_{p\in S_*}p,\quad
 \mathcal D=\{0,1,2401,-343/32\}.
\]
For $a\notin\mathcal D$, let $M_5(a)$ be the product of the
primes $p\notin S_*$ for which at least one of
\[
 v_p(v),\quad v_p(u),\quad v_p(u-v),\quad
 v_p(u-2401v),\quad v_p(32u+343v)
\]
is positive and odd. Every prime occurs only once.
The first valuation is the contact at infinity.

We use the following arithmetic input from
\cite[proof of \texttt{o8-height}, equations
\textup{(LR2)--(LR5)}, and proof of
\texttt{quartic-height19}]{Glazer}.

\begin{definition}[Modular realization input]\label{def:input}
For each $a\in\Q\setminus\mathcal D$, put $R=B_*M_5(a)$.
There is an abelian surface $A$ such that
\begin{equation}\label{eq:geometric}
 h(a)\le2\max(0,\hF(A))+128.
\end{equation}
Either $\hF(A)<854R$, or there are a normal number field $L'/\Q$
and an even nebentypus $\varepsilon$ of a primitive weight-two
newform, of level $N$, such that
\begin{equation}\label{eq:input}
 \begin{gathered}
 [L':\Q]\le2304\varphi(R),\qquad
 N\le210^{300}R^2,\qquad \rad(N)\mid R,\\
 \varepsilon^2|_{G_{L'}}=1,
 \qquad
 A_{L'}\text{ is an isogeny factor of }
       \Jac(X_1(N)/\ker\varepsilon)_{L'}.
 \end{gathered}
\end{equation}
\end{definition}

Here $R$ is squarefree and divisible by $6$.
The degree bound in \eqref{eq:input} retains the totient from the
scalar character: its order satisfies
$m\le2\prod_{p\mid R,\ p\text{ odd}}(p-1)=2\varphi(R)$,
and the construction gives $[L':\Q]\le1152m$.
The square of the nebentypus is trivial over $L'$ by the determinant
identity for the compatible representations. The newform quotient
factors through the diamond quotient because every element of
$\ker\varepsilon$ fixes all conjugates of its differentials.
No bound on the degree of the newform's coefficient field is assumed.

\begin{theorem}\label{thm:contactheight}
Assume the modular realization input. Put
\[
 D_0=4608\cdot210^{300},\qquad
 C_3=19D_0B_*^3\log B_*.
\]
Then $C_3<2^{330770}$ and
\begin{equation}\label{eq:contact3}
 h(a)\le C_3 M_5(a)^3(1+\log M_5(a))
 \quad(a\in\Q\setminus\mathcal D).
\end{equation}
In particular,
\begin{equation}\label{eq:contact-seven-halves}
 h(a)<2^{330789}M_5(a)^{7/2}.
\end{equation}
For every fixed $\alpha>3$ there is an effective constant
$C_\alpha$ such that $h(a)\le C_\alpha M_5(a)^\alpha$.
\end{theorem}
\begin{proof}
Set $M=M_5(a)$ and $R=B_*M$. In the non-CM alternative, let
$H=\ker\varepsilon$ and $t_H=\operatorname{ord}\varepsilon$.
Since $\varepsilon^2$ is trivial on $G_{L'}$, its image order
divides $[L':\Q]$. Hence
\[
 t_H\le2[L':\Q]\le4608\varphi(R).
\]
The exact index identity in Lemma~\ref{lem:modular-model} gives
\begin{equation}\label{eq:cubic-degree}
 \begin{split}
 d_H
 &\le4608\cdot210^{300}R^2\varphi(R)
                    \prod_{p\mid R}(1+p^{-1})\\
 &=D_0R^3\prod_{p\mid R}(1-p^{-2})<D_0R^3.
 \end{split}
\end{equation}
Also $2\le g_H\le d_H$, because the Jacobian contains a surface
factor and the cover is a Belyi map.

Write $\ell=\log R$. The exact prime-product computation in
Section~\ref{sec:verification} gives
\[
 D_0^2<B_*,\qquad144<B_*,\qquad \ell\ge1.
\]
Consequently
\[
 \log N\le\tfrac52\ell,\qquad
 \log d_H\le\tfrac72\ell,\qquad
 \log12\le\tfrac12\ell.
\]
Using $\log(2/\sqrt\pi)<1$ in
Corollary~\ref{cor:modularheight}, we obtain
\begin{equation}\label{eq:jacobian-numeric}
 \hF(J_H)\le5d_H\ell+d_H\le6d_H\ell.
\end{equation}
Thus the number $H_J$ in Lemma~\ref{lem:surface} is at most
\[
 H_*=8D_0R^3\ell.
\]
Let $d'=[L':\Q]$. Since $d'\le2304R$, $2304<e^8$,
$\ell\ge1$, and $\log D_0\le\ell$,
\[
 \log d'<9\ell,\qquad
 \log H_*\le9\ell.
\]
Moreover $H_*>\max(1,\log d')$.
Lemma~\ref{lem:surface}, with $\log14<3$, now gives
\[
 \begin{split}
 \hF(A)
 &\le H_*+8(96+18\ell)\\
 &\le H_*+912\ell
 \le9D_0R^3\ell.
 \end{split}
\]
The CM bound $854R$ is smaller than the same expression.
Combining with \eqref{eq:geometric}, and using
$D_0R^3\ell>128$, gives
\[
 h(a)\le19D_0R^3\ell.
\]
Finally $\log R=\log B_*+\log M\le
(\log B_*)(1+\log M)$, since $\log B_*>1$ and $M\ge1$.
This proves \eqref{eq:contact3}.

The explicit numerical bound uses
\[
 19<2^5,\qquad D_0<2^{2413},\qquad
 B_*<2^{109445},\qquad \log B_*<2^{17}.
\]
Consequently
\[
 C_3<2^{5+2413+3\cdot109445+17}=2^{330770}.
\]
For $u\ge0$, the maximum of $(1+u)e^{-u/2}$ occurs at
$u=1$ and is less than two. Taking $u=\log M$ gives a
constant smaller than $2^{330771}$, and in particular proves
the convenient bound~\eqref{eq:contact-seven-halves}.
More generally, for each $\eta>0$,
$(1+u)e^{-\eta u}$ is bounded on $[0,\infty)$ by an explicitly
computable number. Apply this with $\eta=\alpha-3$.
\end{proof}

\section{Disjoint sets of contacts}

\begin{corollary}\label{cor:amplification}
Fix $r\ge1$ and put
\[
 \Theta_r=\bigcup_{j=0}^{r-1}
 \left(3j+\left\{0,-\tfrac12,-1,
                     \tfrac1{2399},-\tfrac{32}{407}\right\}\right).
\]
There is an explicitly determined finite exceptional set $\Sigma_r$
such that, if $M_r(t)$ is the product of primes outside $\Sigma_r$
at which some $t-\theta$, $\theta\in\Theta_r$, has positive odd
valuation, then
\[
 h(t)\ll_r M_r(t)^{3/r}(1+\log M_r(t))
       \ll_r M_r(t)^{7/(2r)}
       \qquad(t\in\Q\setminus\Theta_r).
\]
All implied constants are effective under the stated modular
realization input.
\end{corollary}
\begin{proof}
Apply Theorem~\ref{thm:contactheight} to
$\phi_j(t)=2+1/(t-3j)$ for $0\le j<r$.
Their matrices have determinant $-1$, so
$h(t)\le h(\phi_j(t))+O_r(1)$.
In addition to $S_*$, exclude the primes dividing contact
coefficients and the nonzero determinants of pairs of primitive
contact forms. The $5r$ rational contacts are distinct.
Outside the resulting fixed set, different contact forms cannot
vanish at the same residue point. Thus the radicals $M_j(t)$
for the different blocks have disjoint prime support, and
$\prod_jM_j(t)=M_r(t)$. Removing the exceptional primes from
a block changes its height inequality by at most an effective
fixed factor.

Multiplying the $r$ inequalities gives
\[
 h(t)^r\ll_r M_r(t)^3
                  \prod_{j=0}^{r-1}(1+\log M_j(t)).
\]
Each logarithmic factor is at most $1+\log M_r(t)$.
Taking $r$th roots proves the first bound. Applying
\eqref{eq:contact-seven-halves} in each block proves the second.
The primitive contact forms, their pairwise determinants, and the
fixed exceptions can be constructed directly as in
\cite[Sections~1--2]{Amplification}.
\end{proof}

Four translated sets of five contacts therefore give height exponent
$7/8$ with twenty distinct finite contacts. This statement concerns
logarithmic height; replacing it by an estimate for multiplicative
height would be a different assertion.

\section{Exact verification}\label{sec:verification}

The script
\begin{center}\small
\nolinkurl{computation/contact_heights/check_evaluation_height.py}
\end{center}
verifies the rational formulas for the discriminant and $j$-invariant
in \eqref{eq:universal-curve}, reconstructs $B_*$ by a sieve with
independent trial division, and checks the integer comparisons
$D_0^2<B_*$, $144<B_*$, $D_0<2^{2413}$, and $B_*<2^{109445}$.
It also checks the exponent arithmetic and the exponential
inequalities used above by rational partial sums.
Its output is
\nolinkurl{evaluation-height-certificate.json} in the same directory.
The finite-algebra and geometric assertions are proved in the text;
they are not inferred from numerical experiments.

\begin{thebibliography}{9}
\bibitem{Stacks}
The Stacks Project Authors, \emph{The Stacks Project},
Tags~00R4, 0AWQ, 0AVT, 0BUL, 0C0I, 0BW0, and 0BWB,
\url{https://stacks.math.columbia.edu}.

\bibitem{Javanpeykar}
A.~Javanpeykar, with an appendix by P.~Bruin,
\emph{Polynomial bounds for Arakelov invariants of Belyi curves},
Algebra \& Number Theory \textbf{8} (2014), 89--140,
\url{https://doi.org/10.2140/ant.2014.8.89}.

\bibitem{Serre}
J.-P.~Serre, \emph{Local Fields}, Graduate Texts in Mathematics
\textbf{67}, Springer, 1979, Chapter~III,
\url{https://doi.org/10.1007/978-1-4757-5673-9}.

\bibitem{vonKaenel}
R.~von K\"anel, \emph{The effective Shafarevich conjecture for abelian
varieties of $\mathrm{GL}_2$-type}, Forum of Mathematics, Sigma
\textbf{9} (2021), e39,
\url{https://doi.org/10.1017/fms.2021.29}.

\bibitem{Remond}
G.~R\'emond, \emph{Propri\'et\'es de la hauteur de Faltings},
Annali della Scuola Normale Superiore di Pisa, Classe di Scienze
\textbf{23} (2022), 1589--1596,
\url{https://doi.org/10.2422/2036-2145.202010_062}.

\bibitem{GaudronRemond}
\'E.~Gaudron and G.~R\'emond,
\emph{Nouveaux th\'eor\`emes d'isog\'enie},
M\'emoires de la Soci\'et\'e Math\'ematique de France
\textbf{176} (2023), Theorem~1.9(1),
\url{https://doi.org/10.24033/msmf.484}.

\bibitem{Glazer}
E.~Glazer and contributors,
\emph{Diophantine definition of the integers in the rationals},
\texttt{manuscripts/quartic\_v9.tex}, commit
\texttt{14e8f601fe0d9173d4a5f084aedc3e07066a7688},
\url{https://github.com/elliotglazer/diophantine-z-q}.

\bibitem{Amplification}
J.~Brox and contributors,
\emph{Height bounds from disjoint sets of rational contacts},
\texttt{contact\_height\_amplification\_20261109\_v13.tex},
\url{https://github.com/JoseBrox/diophantine-z-q_improvements}.
\end{thebibliography}
\end{document}
